% sec-03: outline item C.  Figures 7 (mermaid), 8 (d2), 9 (graphviz).
% The arithmetic spine.  Tables 8 to 11 reconcile exactly.
\section{The \$1.6M Gate and the \$3.5M Programme}

\subsection{One Budget, Reused Without Re-derivation}

The five-year frame is taken verbatim from the funding application of 12 July
2026 and is not re-derived here: \$700,000 per year of direct cost, \$3,500,000
over five years, no cost share, no programme income anticipated
\cite{fundapp2}. It divides into four layers of \$1,600,000 for clinical
conduct, \$720,000 for IND maintenance and safety reporting, \$780,000 for the
interlock rig with its logging and audit replay, and \$400,000 for the
verification package and archive.

\subsection{The Award, and the Work Inside It}

An SBIR award is a total-cost figure and not a work figure. It carries direct
costs, indirect costs, and a fee that is calculated on the sum of the two and is
not itself a cost of the project \cite{nihgpsfee}. A milestone table that
confuses the two will not survive a first reading.

\begin{apptable}
\begin{tabularx}{\textwidth}{>{\raggedright\arraybackslash}X >{\raggedright\arraybackslash}p{2.4cm} >{\raggedright\arraybackslash}p{2.6cm} >{\raggedright\arraybackslash}p{2.4cm}}
\headrow \thead{Component} & \thead{Phase I, 9 months} & \thead{Phase II, 24 months} & \thead{Both, 33 months}\\
\midrule
Direct costs & \$266,000 & \$1,130,000 & \$1,396,000 \\
Indirect at 7.5 percent of direct & \$20,000 & \$85,000 & \$105,000 \\
Subtotal & \$286,000 & \$1,215,000 & \$1,501,000 \\
Fee at 7 percent of subtotal & \$20,000 & \$85,000 & \$105,000 \\
Total award & \$306,000 & \$1,300,000 & \$1,606,000 \\
Overhead and fee load on direct & 15.0 percent & 15.0 percent & 15.0 percent \\
\end{tabularx}
\end{apptable}
\vspace{-0.65cm}
\tabcap{The two awards decomposed. A \$1,606,000 award carries \$1,396,000 of
direct work at a 15.0 percent load, against an SBIR indirect allowance of 40
percent of direct costs that this plan does not claim.}

\begin{appfloat}[!tb]
\begin{appfig}[mermaid-type / state diagram]
% Five states on a horizontal spine at y = 0, two off-spine outcomes at
% y = -2.55.  The two gate exits take 20 and 14 rather than a shared bend,
% because they leave the same node in different directions.
\node[umlinit] (i0) at (-0.55,0) {};
\node[umlstategray,text width=25mm] (u1) at (0.95,0)
  {Unfunded\\0 dollars committed\\2.6 FTE unpaid};
\node[umlstateon,text width=27mm] (u2) at (4.30,0)
  {Phase I active\\306,000 total cost\\months 1 to 9};
\node[umlstate,text width=25mm,line width=1pt] (u3) at (7.55,0)
  {Gate\\four guards evaluated\\no spend authority};
\node[umlstateon,text width=27mm] (u4) at (10.85,0)
  {Phase II active\\1,300,000 total\\months 10 to 33};
\node[umlstatesoft,text width=25mm] (u5) at (14.05,0)
  {Closeout\\1,606,000 expended\\12 artifacts deposited};
\node[umlstatesoft,text width=25mm] (u6) at (6.10,-2.55)
  {Descoped\\resubmit at 3 participants, not 6};
\node[umlstategray,text width=25mm] (u7) at (10.10,-2.55)
  {Terminated\\archive deposited\\no further spend};
\node[umlfinal] (f0) at (12.95,-2.55) {};
\draw[umlarrow] (i0) -- (u1);
\draw[umlarrow] (u1) -- node[umlguard,pos=0.5] {[notice of award]} (u2);
\draw[umlarrow] (u2) -- node[umlguard,pos=0.5] {[month 9]} (u3);
\draw[umlarrow] (u3) -- node[umlguard,pos=0.5] {[G1 G2 G3 G4]} (u4);
\draw[umlarrow] (u4) -- node[umlguard,pos=0.5] {[M12 deposited]} (u5);
\draw[umlarrow] (u3.south west) to[bend right=20]
  node[umlguard,pos=0.4] {[not G4]} (u6.north);
\draw[umlarrow] (u3.south east) to[bend left=14]
  node[umlguard,pos=0.45] {[not G1 or not G3]} (u7.north west);
\draw[umlarrow] (u4.south) -- node[umlguard,pos=0.5] {[DSMB halt 60 days]} (u7.north east);
\draw[umldash] (u6.south) to[bend right=32]
  node[umlguard,pos=0.5] {[site executes, next cycle]} (u3.south);
\draw[umlarrow] (u7) -- (f0);
\node[umlnote,text width=44mm,anchor=north west] at (-0.55,-3.95)
  {G1 bench stop latency at or below 250 ms at p95 over 200 runs. G2 every ASME
   V and V 40 credibility factor at or above gate. G3 IND amendment safe to
   proceed. G4 site agreement executed and IRB approval in hand.};
\node[font=\tiny\itshape,text=protogray,anchor=north west,text width=84mm]
  at (5.10,-3.95) {G4 is the only guard the company cannot satisfy by working
  harder, and the Descoped state exists only because G4 can fail while G1, G2
  and G3 all hold. Every guard label sits above its own transition at a fixed
  offset with a white fill, so no transition line passes through text.};
\end{appfig}
\vspace{-0.65cm}
\figcaption{One award in five states, and the four guards that must all hold\\
at month nine. Two of the four are technical, one is regulatory,\\
and the fourth belongs to an institution the company cannot bind.}
\end{appfloat}

\textbf{G1, bench stop latency.} Two hundred instrumented runs on the interlock
rig must return a 95th-percentile stop latency at or below 250 ms, against the
protocol's arm-stop requirement of 3 ms and system-stop requirement of 500 ms.
The company can satisfy this alone, and M3 pays for it.

\textbf{G2, verification credibility.} Every credibility factor in the
verification suite must sit at or above its ASME V\&V 40 gate, and the suite
must be hashed so the state it passed in is the state that is archived
\cite{asmevv40,ichm15}. The company can satisfy this alone, and M4 pays for it.

\textbf{G3, regulatory.} The IND amendment must reach safe to proceed with the
30-day clock closed and no clinical hold \cite{ecfr2026part312}. The company
files it; the agency decides it.

\textbf{G4, institutional.} The site agreement must be executed and IRB approval
in hand. This is the only guard no amount of company work can satisfy, and it is
why the Descoped state exists at all: G1, G2 and G3 can all hold while G4 fails,
and the correct response to that combination is a smaller Phase II rather than a
termination.

\subsection{Phase I, Nine Months, \$306,000}

\begin{apptable}
\begin{tabularx}{\textwidth}{>{\raggedright\arraybackslash}X >{\raggedright\arraybackslash}p{2.0cm} >{\raggedright\arraybackslash}p{5.4cm}}
\headrow \thead{Line} & \thead{Amount} & \thead{What it buys}\\
\midrule
CEO and sponsor lead, 0.6 FTE, 9 months & \$94,500 & Protocol, IND amendment, site negotiation \\
Verification engineer, 0.5 FTE, 9 months & \$63,000 & Rig build, instrumentation, VVUQ freeze \\
On-premises compute, two GPU nodes & \$38,000 & Simulation and replay capacity, amortized \\
Interlock rig hardware and instrumentation & \$31,000 & The bench that G1 is measured on \\
Regulatory preparation and FDA package & \$27,500 & The amendment G3 turns on \\
Site feasibility, budget, IRB fees & \$24,000 & M1 and M2 \\
Travel and dissemination & \$8,000 & One programme review, one deposit \\
Subtotal, direct & \$266,000 & \\
Indirect at 7.5 percent and fee at 7 percent & \$40,000 & Insurance, audit, accounting, and the fee \\
Total award & \$306,000 & \\
\end{tabularx}
\end{apptable}
\vspace{-0.65cm}
\tabcap{Phase I in eight direct lines plus the load. No participant is consented
and no drug is dispensed inside Phase I: every dollar buys either an instrument,
a verification, or a signature.}

\subsection{Phase II, Twenty-Four Months, \$1,300,000}

\begin{apptable}
\begin{tabularx}{\textwidth}{>{\raggedright\arraybackslash}X >{\raggedright\arraybackslash}p{2.0cm} >{\raggedright\arraybackslash}p{5.4cm}}
\headrow \thead{Line} & \thead{Amount} & \thead{What it buys}\\
\midrule
CEO and sponsor lead, 0.8 FTE, 24 months & \$280,000 & Sponsor obligations, safety reporting \\
Verification engineer, 1.0 FTE, 24 months & \$224,000 & Logging, audit replay, build signing \\
Clinical research coordinator, 0.5 FTE & \$96,000 & Case report forms, deviations, queries \\
Site clinical conduct, 6 participants & \$288,000 & \$48,000 per participant, all visits \\
Pharmacy, specimens, central laboratory & \$104,000 & Dispensing, handling, assay \\
Independent monitoring and DSMB & \$86,000 & Source verification and four DSMB sittings \\
IND maintenance and safety reporting & \$72,000 & Annual report, expedited reports \\
Audit replay, logging and archive build & \$98,000 & M9 and the deposit at M12 \\
Consultants, biostatistics and regulatory & \$34,000 & Fixed fee, not tied to outcome \\
Travel, dissemination and closeout & \$18,000 & Two reviews, one public deposit \\
Subtotal, direct & \$1,130,000 & \\
Indirect at 7.5 percent and fee at 7 percent & \$170,000 & Insurance, audit, accounting, and the fee \\
Total award & \$1,300,000 & \\
\end{tabularx}
\end{apptable}
\vspace{-0.65cm}
\tabcap{Phase II in ten direct lines plus the load. Six participants at \$48,000
each is the largest single line; if the site executes late, Phase II is
resubmitted for three participants, which is the Descoped state of Figure 7.}

\subsection{The Same Work, Priced Twice}

\begin{appfloat}[!tb]
\begin{appfig}[d2-type / layered ledger]
% A three-column ledger read left to right, deliberately unlike Figure 11's
% bottom-to-top stack.  The three money columns take an identical 21mm width
% and 8.4mm height, so a reader comparing a row compares equal areas.
\foreach \i/\y/\lab/\sty in {1/0/{Clinical conduct: site, pharmacy, monitoring}/d2key,
  2/-1.05/{IND maintenance and safety reporting}/d2mid,
  3/-2.10/{Interlock rig, logging, audit replay}/d2soft,
  4/-3.15/{Verification package and archive}/d2gray}{
  \node[\sty,text width=44mm,minimum height=8.4mm,anchor=west] (L\i) at (0,\y) {\lab};}
\begin{scope}[on background layer]
\node[d2cont,fit=(L1)(L4),inner sep=7pt] (lc) {};
\end{scope}
\node[d2title,anchor=south west] at ([yshift=1.2mm]lc.north west) {One programme, four layers};
\foreach \i/\y/\a/\b/\c in {1/0/{\$1,600,000}/{\$612,000}/{\$988,000},
  2/-1.05/{\$720,000}/{\$268,000}/{\$452,000},
  3/-2.10/{\$780,000}/{\$412,000}/{\$368,000},
  4/-3.15/{\$400,000}/{\$104,000}/{\$296,000}}{
  \node[d2cell,minimum width=21mm,text width=18mm,minimum height=8.4mm] at (6.05,\y) {\a};
  \node[d2cell,fill=pablue2,minimum width=21mm,text width=18mm,minimum height=8.4mm] at (8.45,\y) {\b};
  \node[d2cellg,minimum width=21mm,text width=18mm,minimum height=8.4mm] at (10.85,\y) {\c};}
\node[d2title,anchor=south] at (6.05,0.52) {Price A, 5 years};
\node[d2title,anchor=south] at (8.45,0.52) {Price B, SBIR};
\node[d2title2,anchor=south] at (10.85,0.52) {Not reached};
\draw[pagrayd,line width=0.5pt] (4.95,-3.68) -- (11.95,-3.68);
\node[d2cellh,minimum width=21mm,text width=18mm,minimum height=8.4mm] at (6.05,-4.15) {\$3,500,000};
\node[d2cellk,minimum width=21mm,text width=18mm,minimum height=8.4mm] at (8.45,-4.15) {\$1,396,000};
\node[d2cell,fill=pagrayd,minimum width=21mm,text width=18mm,minimum height=8.4mm] at (10.85,-4.15) {\$2,104,000};
\node[font=\tiny\sffamily\bfseries,anchor=east] at (4.85,-4.15) {Total, direct};
\node[font=\tiny\sffamily,anchor=east,text width=42mm,align=right] at (4.85,-4.75)
  {60 months against 33; B plus the third column equals A on every row};
\begin{scope}[shift={(12.55,-3.15)}]
  \hbarrow{2.85}{3.20}{protoblue}{}{100 pct, 3.5M}
  \hbarrow{2.30}{1.28}{pablue1}{}{39.9 pct}
  \hbarrow{1.75}{1.92}{pagraym}{}{60.1 pct}
  \node[font=\tiny\sffamily,text=protogray,anchor=north west,text width=30mm] at (0,1.45)
    {Bars share one 3.2 cm scale, so the lower two sum to the upper exactly.};
\end{scope}
\node[font=\tiny\itshape,text=protogray,anchor=north west,text width=132mm]
  at (0,-5.35) {The left column is reused verbatim from the July 2026 funding
  application and is not re-derived. The middle column is what the SBIR award
  reaches in 33 of the 60 months, and the right column is the arithmetic
  remainder, not an estimate.};
\end{appfig}
\vspace{-0.65cm}
\figcaption{One programme, four layers, and two prices for the same work. The\\
SBIR route reaches 1,396,000 of direct effort inside a 1,606,000\\
award; the third column is the 2,104,000 that it does not reach.}
\end{appfloat}

\begin{apptable}
\begin{tabularx}{\textwidth}{>{\raggedright\arraybackslash}X >{\raggedright\arraybackslash}p{2.4cm} >{\raggedright\arraybackslash}p{2.4cm} >{\raggedright\arraybackslash}p{2.4cm}}
\headrow \thead{Layer} & \thead{Price A, 5 years} & \thead{Inside the award} & \thead{Not reached}\\
\midrule
Clinical conduct: site, pharmacy, monitoring & \$1,600,000 & \$612,000 & \$988,000 \\
IND maintenance and safety reporting & \$720,000 & \$268,000 & \$452,000 \\
Interlock rig, logging, audit replay & \$780,000 & \$412,000 & \$368,000 \\
Verification package and archive & \$400,000 & \$104,000 & \$296,000 \\
Total, direct & \$3,500,000 & \$1,396,000 & \$2,104,000 \\
Term & 60 months & 33 months & 27 months \\
\end{tabularx}
\end{apptable}
\vspace{-0.65cm}
\tabcap{The four-layer budget under two prices. Every row reconciles: the middle
column plus the right column equals the left, and the three totals are the sums
of their own columns.}

The delta is not a shortfall to apologise for. It is the price of the mechanism,
and it decomposes into five identifiable things: twenty-seven additional months
of term, participants seven through eighteen, two further dose levels and the
expansion at the recommended Phase 2 dose, the correlative pharmacology, and the
full verification re-run with the permanent archive.

\subsection{What the Money Reaches}

\begin{appfloat}[!tb]
\begin{appfig}[graphviz-type / dependency DAG]
% Left to right.  No rank carries more than four nodes.  Exactly two edges
% cross, WP4 to WP5 and WP8 to WP10, both at a shallow angle in open canvas.
\foreach \i/\x/\y/\lab/\sty in {1/0/0/{WP1 site feasibility and budget}/gvboxs,
  2/0/-1.35/{WP2 protocol final and IRB package}/gvboxs,
  3/0/-2.70/{WP3 interlock rig bench build}/gvboxs,
  4/0/-4.05/{WP4 VVUQ freeze and hash}/gvboxs}{
  \node[\sty,text width=24mm] (w\i) at (\x,\y) {\lab};}
\node[gvboxm,text width=24mm] (w5) at (2.90,-3.38) {WP5 IND amendment and FDA package};
\begin{scope}[on background layer]
\node[gvcluster,fit=(w1)(w2)(w3)(w4)(w5),inner sep=6pt] (cp1) {};
\end{scope}
\node[gvctitle,anchor=south west] at ([yshift=1mm]cp1.north west)
  {Reachable at 306,000, months 1 to 9};
\node[gvcellh,minimum width=20mm,text width=17mm,anchor=south east]
  at ([xshift=-1mm,yshift=1mm]cp1.north east) {\$306,000};
\foreach \i/\x/\y/\lab/\sty in {6/6.20/0/{WP6 site activation and first dosing}/gvbox,
  7/9.05/0/{WP7 first advised robotic Whipple}/gvbox,
  8/11.90/0/{WP8 dose level 1 cleared, DSMB}/gvbox,
  9/6.20/-1.35/{WP9 audit replay demonstration}/gvbox,
  10/9.05/-1.35/{WP10 dose level 2 cleared, n is 6}/gvbox,
  11/11.90/-1.35/{WP11 interim PK, PD and ctDNA}/gvbox}{
  \node[\sty,text width=24mm] (w\i) at (\x,\y) {\lab};}
\node[gvboxk,text width=24mm] (w12) at (9.05,-2.70) {WP12 closeout and public deposit};
\begin{scope}[on background layer]
\node[gvcluster,fit=(w6)(w7)(w8)(w9)(w10)(w11)(w12),inner sep=6pt] (cp2) {};
\end{scope}
\node[gvctitle,anchor=south west] at ([yshift=1mm]cp2.north west)
  {Reachable at 1,606,000, months 10 to 33};
\node[gvcellh,minimum width=20mm,text width=17mm,anchor=south east]
  at ([xshift=-1mm,yshift=1mm]cp2.north east) {\$1,300,000};
\foreach \i/\x/\y/\lab in {13/6.20/-5.55/{WP13 participants 7 to 18, DL3 and DL4},
  14/9.05/-5.55/{WP14 correlative PK, ctDNA, tissue},
  15/11.90/-5.55/{WP15 full VVUQ re-run and file},
  16/6.20/-6.90/{WP16 second engineer, months 34 to 60},
  17/9.05/-6.90/{WP17 archive, hosting, IND to month 60}}{
  \node[gvboxg,text width=24mm] (w\i) at (\x,\y) {\lab};}
\begin{scope}[on background layer]
\node[gvcluster2,fit=(w13)(w14)(w15)(w16)(w17),inner sep=6pt] (cpg) {};
\end{scope}
\node[gvctitle,text=protogray,anchor=south west] at ([yshift=1mm]cpg.north west)
  {Unreachable at 1,606,000, needs 2,104,000 more};
\node[gvcellg,minimum width=20mm,text width=17mm,anchor=south east]
  at ([xshift=-1mm,yshift=1mm]cpg.north east) {\$2,104,000};
\draw[gvedge] (w1) -- (w2);
\draw[gvedge] (w3.east) -- (w5.west);
\draw[gvedge] (w4.east) -- (w5.west);
\draw[gvedge] (w2.east) to[bend left=10] (w6.west);
\draw[gvedge] (w5.east) to[bend right=10] (w6.west);
\draw[gvedge] (w6) -- (w7);
\draw[gvedge] (w7) -- (w8);
\draw[gvedge] (w8.south) to[bend left=18] (w9.east);
\draw[gvedge] (w8.south) to[bend left=12] (w10.north east);
\draw[gvedge] (w10) -- (w11);
\draw[gvedge] (w11.south) to[bend left=16] (w12.east);
\draw[gvedged] (w10.south) -- (w13.north east);
\draw[gvedged] (w13) -- (w14);
\draw[gvedged] (w9.south) to[bend right=12] (w16.north west);
\draw[gvedged] (w12.south) to[bend left=10] (w17.north);
\draw[gvedged] (w13.south) -- (w16.north);
\draw[protoblue,line width=0.8pt] (4.95,0.85) -- (4.95,-4.85);
\node[font=\tiny\sffamily\bfseries,text=protoblue,fill=protowhite,inner sep=1.5pt]
  at (4.95,-4.60) {gate};
\draw[pagrayd,line width=0.8pt] (4.95,-4.55) -- (13.45,-4.55);
\node[font=\tiny\itshape,text=protogray,anchor=north west,text width=132mm]
  at (0,-7.85) {The money runs out one participant after the second dose level
  clears: WP13 depends on WP10, and two of the five unreachable packages depend
  in turn on WP13. WP1 to WP12 carry the identical cost to milestones M1 to M12,
  so this figure, Figure 8 and Figure 13 can be checked against one another.};
\end{appfig}
\vspace{-0.65cm}
\figcaption{Seventeen work packages and every edge between them. Five close\\
at 306,000, seven more at 1,606,000, and five are unreachable\\
at that figure, all five downstream of the sixth participant.}
\end{appfloat}

Two of the five unreachable packages are not about participants at all. WP15,
at \$228,000, is the full verification re-run against post-trial data and the
credibility file that goes with it; without it the credibility score of 81.9
remains a pre-trial claim. WP17, at \$298,000, is the permanent archive, the
replay hosting, and IND maintenance out to month 60. WP17 is the single package
whose removal changes what \S10 is able to claim, which is why it sits in the
unreachable cluster rather than being quietly absorbed into WP12.
